\section{为什么本章必须更严格}

本章较早的版本，正确地看到了框架是被形成出来的，而不是仅仅被遭遇到的。
但它仍然更像一种机制草图。
如果手稿要主张习惯、解释框架与行为机制是历史性生产出来的结构，
那么它就必须说明：高层组织\emph{如何}产生，什么类型的证据支撑这一主张，
以及类比到人类行为时究竟在哪一点开始从直接测量转入解释性推断。

因此，本章采用一个可见的三层架构：
\begin{enumerate}[nosep]
  \item \textbf{生物学奠基：} 当前生物学与神经科学允许我们如何讨论有组织的状态形成与稳定性。
  \item \textbf{物理与系统类比：} 复杂系统与自组织语言在被谨慎使用时，能够贡献什么。
  \item \textbf{数学与解释结构：} 涌现、路径依赖、亚稳态与多尺度组织在一个有纪律模型中意味着什么。
\end{enumerate}

证据政策仍然保持明确。
本章的主要直接来源是
\emph{Metastability demystified: the foundational past, the pragmatic present and the promising future}
（Nature Reviews Neuroscience, 2024, \textbf{medium}）、
\emph{Associative conditioning in gene regulatory network models increases integrative causal emergence}
（Communications Biology, 2025, \textbf{medium}），以及与之相邻的意识实证章节；后者的强来源为在生物学层面严肃对待状态结构、稳定性与转变提供了正当性。
本章也借助 John Holland 的 \emph{Emergence}、\emph{Introduction to the Theory of Complex Systems}
以及偏数学的动力系统教材。
这些来源都不能直接证明本书完整的行为框架模型。
它们的价值在于约束解释，使其在结构上足够严肃，而不是停留在隐喻式即兴发挥。